\section{Concurrencia y sincronización}\label{sec:concurrencia}

La concurrencia del sistema se organiza alrededor de una separación estricta entre el contexto que detecta un fallo, el que lo resuelve y el que atiende el tráfico de red. Cada grupo de hilos tiene un conjunto acotado de operaciones permitidas y un lugar definido donde puede bloquearse, y esa asignación es la que permite que un acceso inválido de la aplicación se convierta en una transacción de página sin que el manejador de señal tome un cerrojo ni abra un socket. Un nodo en operación ejecuta, con la configuración predeterminada, dos despachadores, cuatro trabajadores, un listener, un hilo de latidos y, cuando la redundancia está activa y una baja lo exige, un hilo de reconstrucción, junto con los hilos de la aplicación.

\subsection{Hilo de señal, despachador y trabajadores}

El manejador de \texttt{SIGSEGV} es la única pieza que corre en contexto de señal y su contrato es deliberadamente mínimo. \texttt{dsm\_sig\_install} lo registra con \texttt{sigaction} y las banderas \texttt{SA\_SIGINFO}, \texttt{SA\_ONSTACK} y \texttt{SA\_RESTART}, y cada hilo que toca la vista reserva una pila alternativa de 64 KiB con \texttt{dsm\_thread\_attach}. El manejador lee \texttt{si\_addr}, comprueba que la dirección caiga dentro de la vista, deduce el tipo de acceso del bit 1 de \texttt{REG\_ERR}, calcula el número de página y llena la ranura TLS del hilo, declarada con modelo \texttt{initial-exec} para que su dirección se resuelva al cargar la biblioteca y no en la primera referencia. Después empuja la ranura a la cola de fallos y se duerme en un bucle de \texttt{futex} sobre el campo \texttt{done}. No hay \texttt{malloc}, ni \texttt{printf}, ni mutex, ni \texttt{mprotect}, ni llamadas de red, y esa restricción no es estilística: \texttt{malloc} puede tomar un cerrojo interno que otro hilo tenía cuando falló, y una llamada de red con plazo puede dormir al manejador sobre una pila de señal que no debe crecer.

La cola de fallos es el único punto de contacto entre el manejador y los despachadores, y es una cola MPMC acotada de 4096 celdas sin cerrojos. Cada celda lleva un contador de secuencia y las posiciones de cabeza y cola se reservan con compare-and-swap, de modo que varios hilos de aplicación pueden fallar a la vez y varios despachadores pueden consumir sin serializarse. El productor avisa con \texttt{sem\_post}, operación async-signal-safe, y el consumidor espera en ese contador con un plazo de 100 ms que le permite revisar el apagado. Si la cola está llena, el manejador termina el proceso con \texttt{\_exit(128 + SIGSEGV)} en lugar de bloquearse dentro de la señal, porque una espera sin salida en ese contexto congelaría al hilo sin que ningún despachador pudiera atenderlo.

Los despachadores son hilos creados por \texttt{dispatcher\_start} en cantidad \texttt{DISPATCHER\_THREADS}, dos por defecto. Cada uno extrae una ranura y llama a \texttt{msi\_local\_fault}; si la página está ocupada, no la rechaza al azar sino que la encola en la cola FIFO de esa página y espera en la condvar con el plazo absoluto de \texttt{LOCK\_TIMEOUT\_MS}, con un tope de 64 reintentos como red de seguridad. Cuando la transacción termina, la finalización incrementa la generación de la página y difunde la condvar, y el despachador reintenta una sola vez sobre el estado ya publicado. Al salir, publica el resultado en la ranura con una escritura de liberación, despierta al hilo con \texttt{FUTEX\_WAKE} y contabiliza la latencia. Un error terminal, como un dueño muerto en \texttt{M}, una ranura inválida o falta de espacio, se escribe en la ranura y el manejador reenvía la señal en lugar de dejar al hilo fallar en bucle.

\begin{figure}[H]
\centering
\begin{tikzpicture}[
    scale=0.78, every node/.style={transform shape},
    >={Latex},
    stage/.style={draw=borderblue, fill=softblue, thick, rounded corners, align=center, font=\footnotesize, minimum width=3.3cm, minimum height=0.85cm},
    kern/.style={draw=borderorange, fill=softorange, thick, rounded corners, align=center, font=\footnotesize, minimum width=3.3cm, minimum height=0.85cm},
    net/.style={draw=bordergreen, fill=softgreen, thick, rounded corners, align=center, font=\footnotesize, minimum width=3.3cm, minimum height=0.85cm},
    flow/.style={->, thick, draw=black!65, font=\tiny, align=center},
    back/.style={->, thick, dashed, draw=black!55, font=\tiny, align=center}
]

\node[stage] (app) at (0, 0) {Acceso de la\\aplicación};
\node[kern] (sig) at (4.3, 0) {Kernel entrega\\\texttt{SIGSEGV}};
\node[kern] (hdl) at (8.6, 0) {Manejador\\\texttt{si\_addr}};
\node[stage] (fq) at (12.9, 0) {Cola de fallos\\sin cerrojos};

\node[stage] (disp) at (12.9, -2.4) {Despachador\\\texttt{msi\_local\_fault}};
\node[net] (home) at (8.6, -2.4) {Home remoto\\worker};
\node[net] (apply) at (4.3, -2.4) {\texttt{page\_apply}\\y \texttt{mprotect}};
\node[kern] (wake) at (0, -2.4) {\texttt{futex\_wake}\\y reintento};

\draw[flow] (app) -- node[above, font=\tiny] {acceso} (sig);
\draw[flow] (sig) -- node[above, font=\tiny] {señal} (hdl);
\draw[flow] (hdl) -- node[above, font=\tiny] {\texttt{push}} (fq);
\draw[flow] (fq) -- node[right, font=\tiny] {\texttt{pop}} (disp);
\draw[flow] (disp) -- node[above, font=\tiny] {RPC} (home);
\draw[flow] (home) -- node[above, font=\tiny] {datos} (apply);
\draw[flow] (apply) -- node[above, font=\tiny] {permiso} (wake);
\draw[back] (wake.north) to[bend right=18] node[midway, left, font=\tiny] {reintento} (app.south);

\end{tikzpicture}
\caption{Recorrido de un fallo de página desde el acceso de la aplicación hasta el worker remoto y el reintento de la instrucción.}
\label{fig:recorrido-fallo}
\end{figure}

Los trabajadores forman un grupo distinto, creado por \texttt{worker\_start} con \texttt{WORKER\_THREADS} hilos, cuatro por defecto, que consumen la cola de mensajes de red. La cola es una lista enlazada protegida por un mutex que solo cubre la manipulación de punteros, con dos semáforos contadores: uno de elementos disponibles y otro de capacidad libre. \texttt{worker\_enqueue} reserva una celda de capacidad con \texttt{sem\_trywait} antes de encolar, de modo que la cota de 4096 mensajes se respeta sin bloquear al listener, y transfiere al trabajador la propiedad del búfer de carga. El trabajador retira el elemento, libera la capacidad y procesa el mensaje fuera del mutex; cada tipo se despacha a su manejador de coherencia, de paridad, de cerrojos o de asignación. Ningún trabajador espera indefinidamente a otro nodo: las RPC que un manejador necesita tienen plazo finito, y los manejadores que solo responden, como una invalidación o una degradación, envían la respuesta y vuelven a la cola. La recuperación de una baja se ejecuta como un mensaje sintético \texttt{NODE\_DOWN} en esta misma cola, porque sus RPC por página no pueden correr en el listener, que es el único que completa las respuestas, ni en el hilo de latidos, cuyo retraso declararía muertos a los demás nodos. La figura~\ref{fig:modelo-hilos} resume ese reparto y las comunicaciones que lo sostienen.

\begin{figure}[H]
\centering
\begin{tikzpicture}[
    scale=0.82, every node/.style={transform shape},
    >={Latex},
    app/.style={draw=borderblue, fill=softblue, thick, rounded corners, align=center, font=\footnotesize, minimum width=3.4cm, minimum height=0.9cm},
    sig/.style={draw=borderorange, fill=softorange, thick, rounded corners, align=center, font=\footnotesize, minimum width=3.4cm, minimum height=0.9cm},
    queue/.style={draw=purple!70!black, fill=softpurple, thick, rounded corners, align=center, font=\footnotesize, minimum width=3.4cm, minimum height=0.9cm},
    worker/.style={draw=bordergreen, fill=softgreen, thick, rounded corners, align=center, font=\footnotesize, minimum width=3.4cm, minimum height=0.9cm},
    flow/.style={->, thick, draw=black!65, font=\scriptsize, align=center},
    back/.style={->, thick, dashed, draw=black!55, font=\scriptsize, align=center}
]

\node[app] (app) at (0, 0) {Hilos de aplicación\\pila alternativa};
\node[sig] (sig) at (4.7, 0) {Manejador de\\\texttt{SIGSEGV}};
\node[queue] (fq) at (9.4, 0) {Cola de fallos\\4096 celdas};
\node[app] (disp) at (14.1, 0) {Despachadores (2)\\\texttt{msi\_local\_fault}};

\node[queue] (wait) at (14.1, -2.7) {Cola de espera\\y turno por página};
\node[worker] (apply) at (9.4, -2.7) {\texttt{page\_apply}\\y \texttt{mprotect}};
\node[worker] (wrk) at (4.7, -2.7) {Trabajadores (4)\\cola de mensajes};
\node[sig] (lst) at (0, -2.7) {Listener\\trama por conexión};

\node[queue] (hb) at (0.0, -4.9) {Latidos y\\membresía};
\node[queue] (rb) at (4.7, -4.9) {Reconstrucción\\de franjas};

\draw[flow] (app) -- node[above, font=\tiny] {fallo} (sig);
\draw[flow] (sig) -- node[above, font=\tiny] {\texttt{push}} (fq);
\draw[flow] (fq) -- node[above, font=\tiny] {\texttt{pop}} (disp);
\draw[flow] (disp) -- node[right, font=\tiny] {\texttt{BUSY}} (wait);
\draw[flow] (disp) -- node[above, font=\tiny] {resuelve} (apply);
\draw[flow] (wrk) -- node[above, font=\tiny] {home} (apply);
\draw[flow] (lst) -- node[above, font=\tiny] {tramas} (wrk);
\draw[back] (apply.south) -- (9.4, -3.7) -- (0, -3.7) node[midway, above, font=\tiny] {\texttt{futex} y reintento} -- (app.south);
\draw[flow] (hb) -- (lst);
\draw[flow] (rb) -- (wrk);

\end{tikzpicture}
\caption{Modelo de hilos de un nodo. El manejador solo registra el fallo, los despachadores lo resuelven, los trabajadores atienden la red y el listener completa las respuestas.}
\label{fig:modelo-hilos}
\end{figure}

\subsection{Protección de estructuras de control}

Las estructuras compartidas se protegen con la primitiva más simple que cubra su patrón de acceso, y ninguna de esas primitivas se conserva mientras se espera a otro nodo. El mutex de exclusión mutua se usa para secciones cortas sobre índices, contadores y colas: la tabla de solicitudes pendientes de red, la lista de trabajadores, la tabla del directorio, la cola de transacciones, la tabla de cerrojos y el índice de franjas. El \texttt{pthread\_rwlock} de la membresía permite que casi todas las operaciones lean la tabla de ranuras en paralelo mientras \texttt{JOIN}, \texttt{LEAVE} y \texttt{NODE\_DOWN} la escriben en exclusiva. Los semáforos POSIX contadores resuelven las esperas con plazo, porque \texttt{sem\_timedwait} devuelve el control al vencer y no exige una condición asociada. Las variables de condición se reservan para las colas de espera por página y por turno, donde varios hilos aguardan un cambio de generación o la liberalización de un turno. El \texttt{futex} queda para el único caso donde un mutex no es válido: la espera del hilo de aplicación dentro del manejador de señal.

Las páginas se protegen con franjas de cerrojo y no con un mutex por página. La vista de 4 GiB tiene más de un millón de páginas, de modo que un arreglo de mutex por página costaría decenas de mebibytes solo en sincronización, mientras que las 256 franjas predeterminadas cubren el mismo estado con una fracción de esa memoria. La página $p$ usa la franja $p \bmod 256$, y esa franja protege el estado local, \texttt{tseq\_seen}, los permisos y el cambio de \texttt{mprotect} de la vista. Un hilo toma una sola franja a la vez y la suelta antes de cualquier llamada de red, de modo que dos páginas distintas pueden avanzar en paralelo aunque compartan franja, y una página concreta serializa sus transiciones locales sin serializar las demás.

La membresía y el anillo tienen su propio orden. \texttt{ring\_home\_of} toma el mutex global del anillo, luego el \texttt{rwlock} de lectura de la membresía, compara la firma de la tabla con la que se usó para construir el arreglo y, si no coincide, toma el \texttt{rwlock} de escritura y reconstruye. La firma resume época, ranura, páginas aportadas, estado y UUID de cada nodo vivo, de modo que un cambio de memoria o una baja invalidan el anillo aunque la época no se haya incrementado. Ninguna ruta toma primero el \texttt{rwlock} y después el mutex del anillo, así que el orden global es único y no hay ciclo entre ambos.

Cada módulo conserva su propio mutex para su estado privado: \texttt{sync.mu} cubre la tabla de cerrojos, la cola de cada cerrojo, la tabla de esperas locales, la tabla de propiedades por hilo y las solicitudes remotas en vuelo; \texttt{raid.mu} cubre la tabla de franjas, la versión, la validez del marco de paridad y la tabla de solicitudes pendientes; \texttt{rebuild\_mu} cubre la cola de ranuras por reconstruir y el estado del hilo. Los contadores de estadística y las banderas de ciclo de vida son atómicos porque se leen y escriben desde varios hilos sin una sección crítica que los agrupe: \texttt{running}, los contadores de fallos, mensajes y bytes, y la época de sincronización que consultan las rutas de envío. La bandera \texttt{busy} del directorio y la secuencia \texttt{tseq} no son atómicas porque viven siempre bajo \texttt{dir\_mu} y bajo la franja de la página, respectivamente.

\subsection{Cola de fallos y cola de espera}

La cola de fallos tiene una capacidad fija y un contrato de progreso: el productor nunca se bloquea y el consumidor nunca depende de que un hilo de aplicación avance. Los índices son contadores de 32 bits que crecen sin reiniciar y se reducen módulo la capacidad, y cada celda publica su disponibilidad con un contador de secuencia, de modo que no existe una comparación de índices que pueda confundir lleno con vacío. El manejador empuja y se duerme en su propia ranura; el despachador extrae, resuelve y despierta esa ranura concreta. La cola no guarda prioridad ni orden distinto al de llegada, y los contadores \texttt{slots\_enqueued} y \texttt{slots\_woken} permiten comprobar que cada fallo encolado termina despertando.

La cola de espera por página resuelve el caso de dos despachadores que encuentran la misma página ocupada. Cada espera es un registro en la pila del despachador, enlazado en una lista FIFO protegida por un único mutex y una variable de condición por página. El despachador captura la generación de la página antes de intentar el fallo, encola su registro y comprueba la generación bajo el mutex; si la finalización ocurrió entre el rechazo y el encolado, la generación ya cambió y la espera retorna de inmediato en lugar de perder el aviso. Al vencer el plazo o al apagarse el nodo, el registro se retira de la lista, de modo que no queda memoria que liberar ni una entrada que pueda despertar a un hilo ya ido. La variable de condición se difunde en cada finalización, y todos los esperadores reintentan una sola vez sobre el estado publicado.

El custodio de una página mantiene una segunda cola, la de turnos, que serializa las transacciones que modifican el directorio. La cola distingue dos clases de solicitante y las ordena por un número de secuencia común: los hilos locales que esperan dentro de una llamada y las solicitudes remotas que llegaron por la red. Cuando el turno está ocupado, una solicitud remota no ocupa un trabajador: se copia el mensaje, se guarda en la lista de diferidos y el trabajador vuelve a la cola de red. Al liberarse el turno, \texttt{home\_txn\_release} elige al siguiente por secuencia, le concede \texttt{busy} y, si es remoto, reencola el mensaje con una bandera interna que indica que ya posee el turno. Esa deferencia es la que evita que cuatro trabajadores queden retenidos esperando una página mientras las respuestas que esa página necesita aguardan en la misma cola. La bandera \texttt{busy} se libera en todas las salidas, incluidas las de error, y el turno se concede antes de cualquier envío de respuesta.

La secuencia de transacción y el turno cumplen funciones distintas. El turno decide quién puede modificar el directorio en un instante; \texttt{tseq} decide qué contenido es más nuevo. \texttt{dir\_tseq\_next} incrementa el contador de la página bajo \texttt{dir\_mu} y salta el cero, que queda reservado para ``sin secuencia''. \texttt{tseq\_accept} compara el valor recibido con \texttt{tseq\_seen} bajo la franja de la página y publica el nuevo valor antes de que el llamador instale los datos; si la instalación falla después, la secuencia ya avanzó y una respuesta vieja no puede colarse en un reintento. La capa de red consulta \texttt{msi\_tseq\_stale} antes de despertar una RPC, de modo que una respuesta duplicada o reordenada no activa un hilo que ya había agotado su plazo. El turno y la secuencia se complementan: el primero impide que dos transacciones observen el mismo estado inicial, y la segunda impide que una respuesta atrasada sustituya un contenido más nuevo.

\subsection{Cerrojos distribuidos}

Un cerrojo distribuido tiene un custodio determinista: \texttt{ring\_home\_of} aplica el anillo de dispersión al identificador del cerrojo, y todos los nodos con la misma membresía obtienen la misma ranura sin coordinación adicional. El custodio mantiene el dueño, la solicitud con la que se concedió y una cola FIFO de 64 entradas con la ranura y el \texttt{req\_id} de cada solicitante. La propiedad local se registra por hilo en una tabla de 256 entradas con \texttt{pthread\_equal}, de modo que dos hilos del mismo proceso compiten igual que dos nodos y solo el hilo que tomó el cerrojo puede liberarlo.

En el camino local, \texttt{sync\_lock} registra una espera en la tabla de 128 entradas, cada una con su propio semáforo, y llama a \texttt{sync\_lock\_local\_request}. Si el cerrojo está libre, no hay migración activa y la cola está vacía, la concesión se registra en el acto y la espera se completa sin pasar por la red. Si está ocupado, la solicitud entra en la cola y el hilo espera con \texttt{sem\_timedwait} durante \texttt{LOCK\_TIMEOUT\_MS}. Al vencer, \texttt{sync\_lock\_local\_withdraw} retira la solicitud de la cola o libera la concesión si ya había sido otorgada, y \texttt{sync\_lock\_handoff} atiende al siguiente. La tabla de esperas se reutiliza y el semáforo se inicializa una sola vez, de modo que una concesión inmediata que dejó un permiso sobrante se limpia antes de reutilizar la entrada.

En el camino remoto, la espera no vive en la tabla del módulo sino en una \texttt{pending\_req\_t} de la tabla de red. El solicitante envía \texttt{LOCK\_REQ} y espera en el semáforo de esa entrada; \texttt{LOCK\_GRANT} se clasifica como respuesta, de modo que el listener la correlaciona por \texttt{req\_id} y despierta al hilo sin encolarla para un trabajador. Si el plazo vence, el hilo retira su entrada de la tabla de red, mira el semáforo una última vez y, si la concesión ya había llegado, la registra como propiedad; si no, envía \texttt{LOCK\_REL} al custodio para que retire la solicitud o libere una concesión vencida. Esa secuencia cubre la carrera entre el vencimiento y la respuesta sin dejar el cerrojo retenido por un solicitante que ya se rindió.

La entrega de una concesión usa acuse y retransmisión acotada. \texttt{sync\_deliver\_grant} registra una espera de acuse, envía \texttt{LOCK\_GRANT} y aguarda un \texttt{LOCK\_REL} con la bandera \texttt{PROBE}; si no llega antes del plazo, reenvía hasta \texttt{REQUEST\_RETRIES} veces sin superar \texttt{LOCK\_TIMEOUT\_MS}. Si no logra entregar, deshace la asignación y pasa al siguiente de la cola, de modo que una petición que nunca llegó no deja el cerrojo retenido. Un \texttt{GRANT} duplicado encuentra la propiedad registrada por \texttt{req\_id} y solo se acusa, sin liberar un cerrojo que el nodo todavía sostiene; un \texttt{GRANT} sin propiedad ni solicitud pendiente se acusa y se devuelve con \texttt{LOCK\_REL} para que el custodio lo reasigne.

La caída del custodio se resuelve con una migración por época. \texttt{sync\_on\_node\_down} deriva una época nueva de la membresía y de la última migración conocida, difunde \texttt{LOCK\_MIGRATE} con la ranura caída y marca esa ranura como muerta en la tabla local antes de calcular custodios, para que el anillo produzca el mismo resultado en todos los nodos sin depender del orden de llegada del aviso. Cada nodo vivo responde con \texttt{LOCK\_STATE} a los custodios que el anillo actualizado determine, con su propiedad local, sus solicitudes remotas en vuelo y sus esperas locales. El detector reúne los informes durante un plazo de un segundo, o hasta que lleguen todos, y \texttt{sync\_migrate\_finalize} libera los cerrojos cuyo dueño era alguno de los nodos caídos, conserva la propiedad y la cola de los demás y reanuda las entregas. Una concesión emitida con una época anterior se descarta en \texttt{sync\_epoch\_stale}, de modo que el custodio caído no puede reescribir el estado reconstruido.

\subsection{Barreras}

La barrera usa el anillo con un dominio separado, \texttt{barrier\_id} $\oplus$ \texttt{0x9E3779B9}, para que un cerrojo y una barrera con el mismo identificador no compartan custodio. El custodio mantiene 64 ranuras con direccionamiento abierto, umbral, contador, mapa de arribos y generación. El primer arribo fija el umbral de la ronda y los siguientes deben anunciar el mismo valor; un participante que anuncie otro número no altera el estado y recibe un \texttt{NACK} con \texttt{DSM\_ST\_INVALID}, que el hilo local consume como error y el participante remoto recibe por la ruta de respuestas. El mapa de bits evita contar dos veces al mismo participante y permite reconstruir la lista de quienes deben ser liberados.

Cuando el contador alcanza el umbral, el custodio cierra la ronda bajo \texttt{sync.mu}: copia la lista de participantes, reinicia el contador y el umbral, incrementa la generación y suelta el mutex antes de enviar nada. Después entrega \texttt{BARRIER\_GO} a cada participante con acuse y retransmisión, y despierta su propia espera con la generación siguiente. El receptor acusa el aviso y llama a \texttt{sync\_pending\_wake\_barrier} con la generación recibida más uno; esa función libera solo las esperas cuya generación sea anterior, de modo que un aviso duplicado o atrasado no destapa una ronda futura. La generación es, por tanto, el equivalente del \texttt{tseq} en las barreras: identifica la ronda cerrada y separa una liberación legítima de una notificación vieja.

La barrera comparte con los cerrojos la tabla de esperas locales y el plazo \texttt{LOCK\_TIMEOUT\_MS}, pero no la tabla de propiedades, porque una barrera no se posee ni se libera. El custodio participa como un participante más: su espera se completa en la misma difusión, sin un mensaje de red hacia sí mismo. Si el custodio cambia por una baja, la barrera no migra su umbral ni su generación: el nuevo custodio abre la ranura con umbral cero y generación uno, y los participantes que lleguen después fijan la ronda nueva. Esa decisión mantiene el mecanismo simple y deja la ronda en curso sujeta a su plazo, sin prometer una continuidad que el estado no conserva.

\subsection{\texttt{dsm\_sync} y persistencia de escrituras}

\texttt{dsm\_sync} recorre la tabla local de páginas y persiste únicamente las que están en estado \texttt{M}. Para cada página copia el estado bajo la franja y suelta la franja antes de escribir, porque la persistencia toma el turno del directorio y puede consultar la red, dos operaciones que no deben ocurrir con la franja tomada. Una página en \texttt{S} o en \texttt{I} no se toca, y el resultado de la llamada refleja el primer fallo de persistencia encontrado: un respaldo incompleto no se reporta como éxito.

Cuando el nodo es custodio de la página, \texttt{msi\_writeback\_page} toma el turno del directorio con \texttt{home\_txn\_acquire}. Si el turno está ocupado, espera en la cola de turnos con plazo en lugar de fallar de inmediato. Con el turno en mano, \texttt{home\_persist\_delta} lee el marco de respaldo anterior, escribe el contenido nuevo y, si la redundancia está activa y el contenido cambió, calcula el delta XOR y llama a \texttt{raid\_on\_write}. Solo después libera el turno y difunde la finalización de la página. La lectura del respaldo, la escritura y el delta forman una sola sección lógica: si dos escrituras de la misma página se intercalaran, el delta se calcularía sobre una versión que ya no es la almacenada y la paridad dejaría de representar el contenido.

Cuando el nodo no es custodio, envía \texttt{WRITEBACK} con la página completa, la secuencia observada y la bandera de datos, y espera el acuse con el plazo de \texttt{REQUEST\_TIMEOUT\_MS}. El custodio aplica el mismo procedimiento: si el turno está libre lo toma, persiste y calcula el delta; si está ocupado, deja el mensaje diferido y el trabajador queda libre, de modo que una sincronización no consume el conjunto de trabajadores esperando una página ajena. Si el receptor ya no es el custodio, responde \texttt{DSM\_ST\_NOT\_HOME} y el emisor propaga el error. La secuencia observada viaja en la cabecera y se usa para sembrar el contador del custodio en un traspaso de contenido, donde la operación se marca como transferencia y se omite el delta porque el contenido no cambió.

La persistencia es el punto donde la coherencia y la redundancia se encuentran. La paridad se calcula sobre la versión persistida y no sobre cada almacenamiento de la aplicación, porque una página puede recibir miles de escrituras antes de una sincronización y un delta por cada una multiplicaría el tráfico sin mejorar la durabilidad. \texttt{dsm\_sync} convierte esa decisión en una frontera explícita: lo que se sincroniza queda en el respaldo y en la paridad, y lo que no se sincroniza sigue siendo una escritura en \texttt{M} cuya única copia es el nodo que la posee. La llamada no toma un cerrojo global ni paraleliza las páginas: recorre la tabla en orden y serializa por página a través del turno del custodio, de modo que dos sincronizaciones concurrentes sobre páginas distintas avanzan sin estorbarse.

\subsection{Prevención de interbloqueos}

El sistema evita los interbloqueos con cuatro reglas que se aplican en todo el código y no dependen de detectar un ciclo después de que ocurra. La primera es que ningún cerrojo local se conserva durante una espera de red. \texttt{msi\_local\_fault} toma la franja, fija el estado pendiente y la suelta antes de enviar la RPC; \texttt{msi\_home\_handle\_request} copia la entrada de directorio bajo \texttt{dir\_mu}, suelta el mutex, hace las RPC de degradación e invalidación y solo después vuelve a tomar \texttt{dir\_mu} para publicar el resultado. \texttt{sync\_lock} suelta \texttt{sync.mu} antes de esperar en su semáforo, la entrega de una concesión no retiene el mutex mientras espera el acuse, y la barrera suelta \texttt{sync.mu} antes de difundir los avisos.

La segunda regla es un orden de adquisición único. La cola de turnos toma \texttt{txn\_mu} y después \texttt{dir\_mu}, y nunca al revés; el anillo toma \texttt{g\_ring\_mu} y después el \texttt{rwlock} de la membresía; la franja de una página se toma sola y se suelta antes de \texttt{wait\_mu}, de \texttt{dir\_mu}, de \texttt{txn\_mu} o de cualquier envío. Ninguna ruta anida la franja con el directorio ni el directorio con la tabla de cerrojos. Esa disciplina convierte el grafo de espera en un orden parcial acíclico, y la única espera que podría parecer circular, la de un custodio que atiende una solicitud y a la vez pide una invalidación a otro nodo, termina en un trabajador remoto que responde sin esperar a su vez.

La tercera regla es que toda espera tiene plazo. Las RPC usan \texttt{REQUEST\_TIMEOUT\_MS}, los envíos \texttt{SEND\_TIMEOUT\_MS}, los cerrojos y las barreras \texttt{LOCK\_TIMEOUT\_MS}, las lecturas de socket \texttt{dsm\_read\_full\_deadline}, y las colas de espera usan \texttt{pthread\_cond\_timedwait} o \texttt{sem\_timedwait}. La única espera sin plazo explícito es el \texttt{futex} del hilo de aplicación dentro del manejador, y su progreso no depende de otro hilo de aplicación sino del despachador, que a su vez tiene un tope de 64 reintentos y un plazo absoluto por ranura. Al agotarse, el fallo se declara terminal y el proceso recibe la señal en lugar de quedar bloqueado.

La cuarta regla es que los acuses se resuelven en el listener y no en un trabajador. Un \texttt{LOCK\_REL} con \texttt{PROBE}, un \texttt{BARRIER\_GO} con \texttt{PROBE}, un \texttt{INVALIDATE\_ACK} y una respuesta de sincronización del asignador se atienden en línea, porque si se encolaran, un trabajador podría quedar esperando un acuse que solo otro trabajador podría entregar. La misma razón lleva la migración de cerrojos a un trabajador y no al hilo de latidos: el hilo de latidos no debe esperar a la red, y el trabajador que ejecuta la recuperación no es el mismo que completa las respuestas de página. El manejador de señal queda fuera del grafo por construcción: no toma cerrojos, no hace red y solo comparte con el despachador una ranura y un \texttt{futex}.

\subsection{Interacción entre señales, sockets y memoria}

La convivencia entre señales y red se apoya en un conjunto pequeño de decisiones. \texttt{SIGPIPE} se ignora al inicializar el nodo y cada envío usa \texttt{MSG\_NOSIGNAL}, de modo que un par que cierra a mitad de una transferencia produce \texttt{EPIPE} y no termina el proceso. Las lecturas y escrituras completas reintentan \texttt{EINTR}, y las esperas con plazo se reanudan tras una interrupción en lugar de devolver un error espurio. \texttt{SA\_RESTART} evita que una señal sin manejador interrumpa una llamada al sistema, y el manejador de \texttt{SIGSEGV} se ejecuta sobre la pila alternativa del hilo, que se deshabilita con \texttt{SS\_DISABLE} antes de liberarla en \texttt{dsm\_thread\_detach}.

El orden de \texttt{mprotect} forma parte del protocolo y se aplica en las dos direcciones. Cuando una copia en \texttt{M} se degrada para entregar su contenido, el custodio cierra primero la escritura con \texttt{PROT\_READ}, copia después y solo entonces marca la página como \texttt{S}; si copiara antes de cerrar, un hilo de la aplicación podría escribir entre la copia y el cambio de permiso, y el contenido entregado no incluiría esa escritura. En la dirección contraria, \texttt{page\_apply} instala el contenido en el alias, aplica el permiso nuevo a la vista y publica el estado bajo la misma franja; una página en estado pendiente tiene \texttt{PROT\_NONE}, de modo que un acceso concurrente vuelve a fallar y se encola en lugar de leer un contenido a medias. El alias de escritura conserva \texttt{PROT\_READ | PROT\_WRITE} durante toda la vida del nodo porque es un mapeo distinto del mismo \texttt{memfd}, y el cambio de permisos de la vista no lo afecta.

La correlación por \texttt{req\_id} y la compuerta de secuencia separan las respuestas legítimas de las tardías. El listener valida la cabecera y la longitud de la carga antes de clasificar el mensaje; si es una respuesta, busca la solicitud pendiente y la completa, y si no la encuentra descarta \texttt{PAGE\_DATA} y \texttt{GRANT} para que no lleguen a un trabajador. Una respuesta de página cuya secuencia no supera la última aplicada se rechaza en \texttt{pending\_complete} sin despertar la RPC, y aunque llegara al camino de la página, \texttt{page\_accept\_pending} exige que la página siga en un estado pendiente. Esa doble comprobación impide que una respuesta de una RPC vencida sobrescriba una copia más nueva, y el contador \texttt{tseq\_discarded} deja rastro de cuántas se descartaron.

La memoria de la vista y la de los módulos siguen ciclos de vida distintos que la sincronización debe respetar. \texttt{dsm\_free} no libera una página que está en un estado pendiente: devuelve \texttt{DSM\_ERR\_BUSY} y deja que la aplicación reintente, porque una transacción en curso observaría un contenido a medias. Antes de devolver el rango, la página en \texttt{M} se persiste y, si el nodo es custodio, se invalidan las copias remotas; solo entonces se pone en ceros, con \texttt{PROT\_NONE} y estado \texttt{I}, y el rango vuelve al mapa global. El contexto y sus tablas se liberan después de detener despachadores, trabajadores y listener, en ese orden, porque un hilo de servicio que siga leyendo \texttt{ctx->cfg} durante la destrucción accedería a memoria ya liberada. La cola de fallos se destruye únicamente después de restaurar el manejador anterior, de modo que ningún manejador tardío pueda hacer \texttt{sem\_post} sobre un semáforo destruido.
